% !PRINTED 649
\weberpage{649}
Als Grundlage dient uns dabei die Function
\weq{1}{w = (k_0 - k_2)(k_1 - k_3),}
die von Null verschieden sein muss, und die für einen reellen Körper $\mathfrak{K}$
immer reell ist, da entweder $k_0, k_1, k_2, k_3$ reell sind oder $k_0, k_2$ und
$k_1, k_3$ zwei conjugirt imaginäre Paare bilden.

Bei der cyklischen Permutation $(0, 1, 2, 3)$ ändert $w$ sein Vorzeichen. Also ist das
Quadrat von $w$ eine cyklische Function, die nach Voraussetzung bekannt sein soll. Wir
setzen daher, indem wir durch die kleinen lateinischen Buchstaben $a, b, c \ldots$
Grössen in $\mathfrak{K}$, also rationale Grössen bezeichnen,
\weq{2}{w = 4\sqrt{c},}
und bemerken noch, dass jede Function der $k$, die durch die cyklische Permutation
$(0, 1, 2, 3)$ ihr Vorzeichen ändert, das Product einer rationalen Grösse mit
$\sqrt{c}$ ist.

Es handelt sich dann nur darum, cyklische Functionen in genügender Anzahl zu bilden, so
dass man die vier Grössen $k_0, k_1, k_2, k_3$ daraus berechnen und durch von einander
unabhängige rationale Grössen algebraisch ausdrücken kann.

Nun sind zwei weitere cyklische Functionen
\[
(k_0 - k_2)^2 + (k_1 - k_3)^2, \qquad \frac{(k_0 - k_2)^2 - (k_1 - k_3)^2}{(k_0 - k_2)(k_1 - k_3)},
\]
und wir bekommen also, wenn $a, b$ rationale Grössen sind,
\weq{3}{\begin{aligned}
(k_0 - k_2)^2 + (k_1 - k_3)^2 &= 8\,b \\
(k_0 - k_2)^2 - (k_1 - k_3)^2 &= 8\,a\sqrt{c}.
\end{aligned}}
Die drei rationalen Grössen $a, b, c$ sind aber nicht von einander unabhängig, sondern
es besteht nach (2) und (3) zwischen ihnen die Relation
\weq{4}{b^2 = c\,(1 + a^2).}
Aus (3) ergiebt sich ferner
\weq{5}{\begin{aligned}
k_0 - k_2 &= 2\sqrt{b + a\sqrt{c}} \\
k_1 - k_3 &= 2\sqrt{b - a\sqrt{c}},
\end{aligned}}
und zwischen den beiden Quadratwurzeln besteht die Relation
\weq{6}{\sqrt{c} = \sqrt{b + a\sqrt{c}}\; \sqrt{b - a\sqrt{c}}.}
Bezeichnen wir ferner mit $C$ und $B$ wieder zwei rationale Grössen, so ist
%CONT>>
